\subsection{Education and access to research networks}
Access to research depends on networks as well as money.
Someone can afford a capable computer and still lack a university contact,
a research team, a route to project support or an audience for the work.
Buying hardware does not automatically provide that network. This is why
the university and institutional pathway matters to this census, alongside
the working-class access question.

The Department for Education's July 2026 release follows pupils aged 15
in England's state-funded and special schools in 2020/21 into higher
education by age 19 in 2024/25 \cite{dfe2026}. Overall progression was
44.7\% of 574,876 pupils. Among those recorded as eligible for free school
meals (FSM) at 15, it was 28.6\% of 107,873; among other pupils, it was
48.4\% of 467,003.

The full breakdown in Table~\ref{tab:dfe} shows differences in progression
that the national average hides. Among FSM pupils,
the rates were 54.7\% for Black African boys and 71.8\% for Black African
girls, compared with 26.0\% and 44.2\% for Black Caribbean boys and girls.
The corresponding White British rates were 13.0\% and 21.0\%. Among
non-FSM pupils, the rates were 67.2\% and 81.3\% for Black African boys
and girls, 37.8\% and 58.4\% for Black Caribbean boys and girls, and
35.9\% and 48.6\% for White British boys and girls.

Those differences make detailed categories necessary for interpreting the
access figures. Table~\ref{tab:dfe} therefore includes mixed-background
categories, the detailed Asian and White categories, other ethnic groups
and unknown ethnicity, as well as the published major-group totals. Every
recorded category is visible. The accompanying data retain all 222 original
observations, including the pupil denominators and higher-education entrant
counts behind each percentage.

\begin{table}[htbp]
\begingroup\small\setlength{\tabcolsep}{4pt}
\renewcommand{\arraystretch}{1.16}
\caption{Higher-education progression by age 19 in 2024/25 (\%). England state-funded and special-school pupils aged 15 in 2020/21; every recorded major and detailed ethnicity group. Source: DfE~\cite{dfe2026}.}\label{tab:dfe}
\medskip
\begin{tabular}{@{}>{\raggedright\arraybackslash}p{0.38\linewidth}rrrrrr@{}}
\toprule & \multicolumn{2}{c}{FSM eligible} & \multicolumn{2}{c}{Non-FSM} & \multicolumn{2}{c}{All sexes}\\
Recorded ethnicity & Male & Female & Male & Female & FSM & All\\\midrule
\textbf{White (all)} & \textbf{14.4} & \textbf{22.6} & \textbf{36.9} & \textbf{49.6} & \textbf{18.4} & \textbf{39.0}\\
\hspace*{0.7em}White British & 13.0 & 21.0 & 35.9 & 48.6 & 16.9 & 37.8\\
\hspace*{0.7em}Irish & 20.1 & 26.0 & 52.1 & 66.2 & 23.0 & 53.7\\
\hspace*{0.7em}Traveller of Irish Heritage & 8.6 & 4.8 & 20.9 & 23.3 & 7.0 & 13.3\\
\hspace*{0.7em}Gypsy / Roma & 23.7 & 24.4 & 34.7 & 33.2 & 24.1 & 29.7\\
\hspace*{0.7em}Other White background & 36.0 & 48.9 & 48.1 & 60.7 & 42.4 & 52.8\\
\addlinespace[3pt]
\textbf{Mixed / Multiple ethnic groups (all)} & \textbf{26.1} & \textbf{38.2} & \textbf{48.0} & \textbf{60.7} & \textbf{32.1} & \textbf{48.7}\\
\hspace*{0.7em}White and Black Caribbean & 15.6 & 27.6 & 32.0 & 45.8 & 21.8 & 33.3\\
\hspace*{0.7em}White and Black African & 32.3 & 46.8 & 49.4 & 63.8 & 39.5 & 52.0\\
\hspace*{0.7em}White and Asian & 28.1 & 40.7 & 55.8 & 66.8 & 34.1 & 55.8\\
\hspace*{0.7em}Other Mixed / Multiple background & 32.7 & 44.3 & 52.2 & 65.3 & 38.5 & 53.9\\
\addlinespace[3pt]
\textbf{Asian / Asian British (all)} & \textbf{46.7} & \textbf{61.3} & \textbf{63.4} & \textbf{76.0} & \textbf{53.9} & \textbf{66.4}\\
\hspace*{0.7em}Indian & 46.8 & 61.0 & 69.0 & 80.5 & 53.8 & 72.8\\
\hspace*{0.7em}Pakistani & 42.2 & 58.3 & 53.4 & 68.7 & 50.1 & 58.2\\
\hspace*{0.7em}Bangladeshi & 51.8 & 65.0 & 64.8 & 75.7 & 58.5 & 66.6\\
\hspace*{0.7em}Chinese & 73.4 & 76.4 & 81.8 & 87.8 & 75.2 & 84.1\\
\hspace*{0.7em}Other Asian background & 51.5 & 64.1 & 70.3 & 81.6 & 57.3 & 72.5\\
\addlinespace[3pt]
\textbf{Black / African / Caribbean / Black British (all)} & \textbf{46.3} & \textbf{63.9} & \textbf{60.2} & \textbf{75.8} & \textbf{55.1} & \textbf{64.2}\\
\hspace*{0.7em}Caribbean & 26.0 & 44.2 & 37.8 & 58.4 & 35.1 & 43.7\\
\hspace*{0.7em}African & 54.7 & 71.8 & 67.2 & 81.3 & 63.2 & 71.2\\
\hspace*{0.7em}Other Black / African / Caribbean background & 41.9 & 60.8 & 57.3 & 72.1 & 51.0 & 60.6\\
Any other ethnic group & 49.7 & 62.6 & 57.8 & 71.4 & 56.0 & 61.7\\
Unknown ethnicity & 28.0 & 36.1 & 43.8 & 52.5 & 32.0 & 44.4\\
\addlinespace[3pt]
\textbf{All ethnic groups} & \textbf{23.5} & \textbf{33.9} & \textbf{42.2} & \textbf{54.9} & \textbf{28.6} & \textbf{44.7}\\
\bottomrule\end{tabular}
\par\medskip
\noindent\footnotesize Bold rows are major-group or national totals; indented rows are their detailed categories. Totals and components overlap. White British abbreviates English / Welsh / Scottish / Northern Irish / British. Non-FSM is the source category All Other Pupils. FSM eligibility and sex are recorded at age 15. One pupil with unknown recorded sex is included in the all-sex totals. The released CSVs retain all 222 source rows, original category labels, HE entrant counts and pupil denominators for every displayed percentage.
\par\endgroup
\end{table}


\begin{samepage}
FSM is the published administrative low-income indicator used in these
figures. Working-class disadvantage extends beyond that eligibility
threshold. The June 2026 \emph{Independent Inquiry into White Working Class
Educational Outcomes} examines both White British FSM pupils and the wider
C2DE socioeconomic categories \cite{wwci}. Working-class pupils therefore
appear in both FSM and non-FSM columns. This is an important part of the
access question: people can fall outside a formal eligibility category
while still facing restricted opportunity.\par
\end{samepage}

These educational comparisons provide context for the routes through
which prospective researchers may pass. A capability map confined to
university networks cannot capture the full range of people learning and
building outside them. Finding that talent requires following the work
those people have actually produced. An
independent builder should have a route to support before they have to
buy advertising, establish a large social-media audience or become a
political campaign's example.
